\section{Discussion}
\label{sec:discussion}

\subsection{Hypothesis Evaluation}

All five pre-registered predictions were confirmed:

\begin{enumerate}
\item[\textbf{P1}] Integrated verification with Layout A uses 285,756 gas (first batch)
and 268,656 gas (steady state), both under 300K. \textbf{Confirmed.}

\item[\textbf{P2}] Delegated verification is analytically estimated at $\sim$341,856 gas,
exceeding 300K by 14\%. \textbf{Confirmed} (analytically; not directly benchmarked).

\item[\textbf{P3}] Layout A uses 32 bytes per batch (1 storage slot), well under the
100-byte target. \textbf{Confirmed.}

\item[\textbf{P4}] Layout B uses 64 bytes per batch (2 storage slots), well under the
200-byte target. \textbf{Confirmed.}

\item[\textbf{P5}] The root chain continuity check ($\mathit{prevRoot} =
\mathit{currentRoot}$) catches 100\% of gap and reversal attempts.
\textbf{Confirmed} (7/7 invariant tests passed).
\end{enumerate}

\subsection{The 72\% Verification Dominance}

The most significant finding is the gas budget decomposition: 71.9\% verification, 28.1\%
storage and logic. This has three important architectural implications.

First, \emph{storage layout optimization has diminishing returns}. The difference between
the cheapest layout (Events Only, 57,887 gas storage) and the most expensive (Rich,
102,799 gas storage) is 44,912 gas---but this represents only 15.7\% of the total budget.
Even eliminating all storage operations entirely would save less gas than a single pairing
check (181,000 gas).

Second, \emph{integrated verification is a hard constraint, not a preference}. Delegating
verification to a separate contract adds $\sim$56K gas from the cross-contract call and
redundant storage, pushing even the minimal layout over 300K. There is no storage layout
that compensates for the delegation overhead within a 300K budget.

Third, \emph{future gas optimization must target the verification cost}. Migrating from
Groth16 to a proof system with cheaper on-chain verification (e.g., PLONK with KZG
commitments verified via a single pairing, or recursive proof aggregation that amortizes
verification across batches) would expand the storage budget proportionally.

\subsection{Layout A as the Optimal Trade-off}

Layout A (Minimal) achieves the best trade-off among the three candidates:

\begin{itemize}
\item \textbf{vs. Layout C (Events Only):} Layout A adds 22,269 gas (+8.4\% of total) to
store each batch's state root on-chain. This cost buys \emph{direct on-chain queryability}:
any contract or cross-enterprise verifier can call \texttt{getBatchRoot(enterprise, batchId)}
to verify a historical state root without relying on event log indexing. For an enterprise
system where cross-enterprise verification and audit compliance are requirements, this
on-chain availability is essential.

\item \textbf{vs. Layout B (Rich):} Layout B adds 22,643 gas (+7.9\% of total) over
Layout A to store batch metadata (size, timestamp, cumulative transactions) on-chain.
This metadata is already available via event logs, which are queryable through standard
Ethereum JSON-RPC (\texttt{eth\_getLogs}). The marginal benefit of on-chain metadata
storage does not justify exceeding the 300K target on first batch submissions.
\end{itemize}

\subsection{Single-Phase vs. Multi-Phase Commitment}

Our single-phase atomic design diverges from all three production systems, which use 2--3
phases. This is not an oversight but a deliberate architectural choice enabled by the
enterprise validium trust model:

\begin{itemize}
\item \textbf{No sequencer-prover separation:} In public rollups, the sequencer and prover
are typically different entities, necessitating a commit phase where the sequencer posts
data and a prove phase where the prover submits the proof. In our model, the enterprise
node performs both roles atomically.

\item \textbf{No challenge period:} Public rollups may need to delay execution to allow
for fraud proofs or proof verification disputes. Enterprise validium verifies the proof
on-chain in the same transaction, achieving immediate finality.

\item \textbf{Reduced L1 transaction count:} A single L1 transaction per batch (vs. 2--3
for public rollups) halves the per-batch overhead on the L1 chain, which matters even in
a zero-fee model because of block gas limits and throughput constraints.
\end{itemize}

The trade-off is that our model cannot aggregate proofs across multiple committed-but-unproved
batches. Each batch requires its own proof. Future work on recursive proof aggregation would
enable amortizing a single on-chain verification across multiple batches while preserving the
single-transaction submission model.

\subsection{Comparison with Production Systems}

Our per-batch storage of 32 bytes compares favorably with production systems (64--96 bytes),
achieved through the simpler trust model that eliminates the need for batch hashes, status
flags, and withdrawal roots. However, this comparison must be qualified: production systems
store additional data to support features (L2$\to$L1 messaging, withdrawal proofs,
challenge windows) that enterprise validium does not require.

The gas comparison is less direct because production systems amortize verification across
many batches. zkSync Era, for example, can prove hundreds of committed batches in a single
\texttt{proveBatches} call. Our per-batch verification cost of 205,600 gas would be
amortized to $205{,}600 / B$ per batch with $B$-batch aggregation, but this requires
recursive proofs or proof aggregation---a separate research question.

\subsection{Limitations}

\textbf{Mock verification.} ZK verification gas was added analytically rather than measured
end-to-end. While EIP-1108 precompile costs are fixed constants, the Solidity overhead for
proof deserialization and precompile invocation may add 1--3K gas not captured in our model.

\textbf{Single EVM implementation.} All benchmarks were run on Hardhat's JavaScript EVM.
Gas costs on Subnet-EVM (go-ethereum based) should be identical per the specification, but
minor implementation differences in gas accounting are possible.

\textbf{No production calldata.} Our benchmarks used minimal calldata (two 32-byte roots +
one uint256 batch size). A production submission with a full Groth16 proof ($\sim$800 bytes)
adds approximately 12,800 gas in calldata costs (800 nonzero bytes $\times$ 16 gas/byte),
which is accounted for in our analytical budget but not directly measured.

\textbf{Analytical delegated verification.} P2 (delegated verification exceeds 300K) was
confirmed analytically, not by deploying and benchmarking a full \texttt{ZKVerifier}
contract. The estimate is conservative (likely underestimates the actual overhead) because
it does not account for function selector dispatch, argument encoding, and return value
decoding in the cross-contract call.
